\section{Arc Length and Surface Area}\label{sec:arc_length}

In previous sections we have used integration to answer the following questions:
\begin{enumerate}
	\item	Given a region, what is its area?
	\item	Given a solid, what is its volume?
\end{enumerate}

In this section, we address a related question: Given a curve, what is its length? This is often referred to as \textbf{arc length}. 

\mtable{Graphing $y=\sin x$ on $[0,\pi]$ and approximating the curve with line segments.}{fig:arcintro}{%
\begin{tikzpicture}[alt={Blue y = sin x arch on 0 < x < π, peak at (π/2 , 1); ticks at π/4, π/2, 3π/4.}]
 \begin{axis}[width=\marginparwidth,%height=.46\marginparwidth,
   tick label style={font=\scriptsize},axis equal,
   axis y line=middle,axis x line=middle,name=myplot,axis on top,
   xtick={.79,1.57,2.36,3.14},
   xticklabels={$\frac{\pi}4$,$\frac{\pi}2$,$\frac{3\pi}{4}$,$\pi$},
   ymin=-.2,ymax=1.25,
   xmin=-.1,xmax=3.5]
  \addplot [draw={\colorone},smooth,thick,domain=0:3.14] {sin(deg(x))};
 \end{axis}
 \node [right] at (myplot.right of origin) {\scriptsize $x$};
 \node [above] at (myplot.above origin) {\scriptsize $y$};
\end{tikzpicture}
\\(a)\\
\begin{tikzpicture}[alt={Same sin curve; red polyline through five black nodes approximates it on 0 < x < π.}]
 \begin{axis}[width=\marginparwidth,%height=.46\marginparwidth,
   tick label style={font=\scriptsize},
   axis y line=middle,axis x line=middle,name=myplot,axis on top,%
   xtick={.79,1.57,2.36,3.14},
   xticklabels={$\frac{\pi}4$,$\frac{\pi}2$,$\frac{3\pi}{4}$,$\pi$},
   ytick={.71,1},%
   yticklabels={$\frac1{\sqrt2}$,$1$},axis equal,
   ymin=-.2,ymax=1.25,
   xmin=-.1,xmax=3.5]
  \addplot [draw={\colorone},smooth,thick,domain=0:3.14] {sin(deg(x))};
  \draw [draw={\colortwo},thick] (axis cs:0,0) -- (axis cs:.79,.71) -- (axis cs: 1.57,1) -- (axis cs:2.36,.71) -- (axis cs: 3.14,0);
  \filldraw (axis cs:0,0) circle (1pt) (axis cs:.79,.71) circle (1pt) (axis cs: 1.57,1) circle (1pt) (axis cs:2.36,.71) circle (1pt) (axis cs: 3.14,0)circle (1pt);
 \end{axis}
 \node [right] at (myplot.right of origin) {\scriptsize $x$};
 \node [above] at (myplot.above origin) {\scriptsize $y$};
\end{tikzpicture}
\\(b)}

Consider the graph of $y=\sin x$ on $[0,\pi]$ given in \autoref{fig:arcintro} (a). How long is this curve? That is, if we were to use a piece of string to exactly match the shape of this curve, how long would the string be?

As we have done in the past, we start by approximating; later, we will refine our answer using limits to get an exact solution.

The length of straight-line segments is easy to compute using the Distance Formula. We can approximate the length of the given curve by approximating the curve with straight lines and measuring their lengths. 

In \autoref{fig:arcintro} (b), the curve $y=\sin x$ has been approximated with 4 line segments (the interval $[0,\pi]$ has been divided into 4 equal length subintervals). It is clear that these four line segments approximate $y=\sin x$ very well on the first and last subinterval, though not so well in the middle. Regardless, the sum of the lengths of the line segments is $3.79$, so we approximate the arc length of $y=\sin x$ on $[0,\pi]$ to be $3.79$. 

%Using 6 evenly spaced subintervals is not too much more work and gives a better approximation. The six lines are shown in \autoref{fig:arcintro} (c) and provide an arc length approximation of $3.81$. 

In general,  we can approximate the arc length of $y=f(x)$ on $[a,b]$ in the following manner. Let $a=x_0 < x_1 < \dotso < x_{n-1}< x_n=b$ be a partition of $[a,b]$ into $n$ subintervals. Let $\Delta x_i$ represent the length of the $i^{\text{th}}$ subinterval $[x_{i-1},x_i]$.

\autoref{fig:arcintro2} zooms in on the $i^{\text{th}}$ subinterval where $y=f(x)$ is approximated by a straight line segment. The dashed lines show that we can view this line segment as the hypotenuse of a right triangle whose sides have length $\Delta x_i$ and $\Delta y_i$. Using the Pythagorean Theorem, the length of this line segment is $\ds \sqrt{(\Delta x_i)^2 + (\Delta y_i)^2}.$ Summing over all subintervals gives an arc length approximation
%
\mtable{Zooming in on the $i^{\text{th}}$ subinterval $[x_{i-1},x_i]$ of a partition of $[a,b]$.}{fig:arcintro2}{\begin{tikzpicture}[alt={Zoom on x_{i−1} < x < x_i: blue arc, red chord, labeled delta x horizontal and delta y vertical.}]
 \begin{axis}[width=\marginparwidth,tick label style={font=\scriptsize},
   axis y line=middle,axis x line=middle,name=myplot,axis on top,
   xtick=\empty,extra x ticks={.3,1.37},extra x tick labels={$x_{i-1}$,$x_i$},
   ytick=\empty,extra y ticks={.48,1},extra y tick labels={$y_{i-1}$,$y_i$},
   ymin=-.2,ymax=1.25,xmin=-.1,xmax=1.6]
  \addplot [draw={\colorone},smooth,thick,domain=0.3:1.37] {sin(deg(x+.2))};
  \draw (axis cs: .25,1) -- (axis cs:.1,1) 
        (axis cs: .25,.48) -- (axis cs:.1,.48)
        (axis cs: .175,.48)
         -- node [pos=.5,fill=white] {\scriptsize $\Delta y_i$} (axis cs:.175,1);
  \draw (axis cs: .3,.25) -- (axis cs: .3,.4)
        (axis cs: 1.37,.25) -- (axis cs: 1.37,.4)
        (axis cs: .3,.325)
         -- node [pos=.5,fill=white] {\scriptsize $\Delta x_i$} (axis cs: 1.37,.325);
  \draw [draw={\colortwo},thick] (axis cs: .3,.48) -- (axis cs:1.37,1);
  \draw [dashed,thin] (axis cs: .3,.48) -| (axis cs:1.37,1);
  \filldraw (axis cs: .3,.48) circle (1pt) (axis cs:1.37,1) circle (1pt);
 \end{axis}
 \node [right] at (myplot.right of origin) {\scriptsize $x$};
 \node [above] at (myplot.above origin) {\scriptsize $y$};
\end{tikzpicture}}
%
\[L \approx \sum_{i=1}^n \sqrt{(\Delta x_i)^2 + (\Delta y_i)^2}.\]

As it is written, this is \emph{not} a Riemann Sum. While we could conclude that taking a limit as the subinterval length goes to zero gives the exact arc length, we would not be able to compute the answer with a definite integral. We need first to do a little algebra.

In the above expression factor out a $\Delta x_i^2$ term:
\begin{align*}
	\sum_{i=1}^n \sqrt{(\Delta x_i)^2 + (\Delta y_i)^2}
	&= \sum_{i=1}^n \sqrt{(\Delta x_i)^2\left(1 + \frac{(\Delta y_i)^2}{(\Delta x_i)^2}\right)}.
\end{align*}
Now pull the $(\Delta x_i)^2$ term out of the square root:
\[L\approx\sum_{i=1}^n\sqrt{1 + \frac{(\Delta y_i)^2}{(\Delta x_i)^2}}\ \Delta x_i.\]
This is nearly a Riemann Sum. Consider the $(\Delta y_i)^2/(\Delta x_i)^2$ term. The expression $\Delta y_i/\Delta x_i$ measures the ``change in $y$/change in $x$,'' that is, the ``rise over run'' of $f$ on the $i^{\text{th}}$ subinterval. The Mean Value Theorem of Differentiation (\autoref{thm:mvt}) states that there is a $c_i$ in the $i^{\text{th}}$ subinterval where $\fp(c_i) = \Delta y_i/\Delta x_i$. Thus we can rewrite our above expression as:
\[L\approx\sum_{i=1}^n\sqrt{1+[\fp(c_i)]^2}\ \Delta x_i.\]
This \emph{is} a Riemann Sum. As long as \fp\ is continuous on $[a,b]$, we can invoke \autoref{thm:riemannSum} and conclude
\[L=\int_a^b\sqrt{1+[\fp(x)]^2}\dd x.\]

\begin{keyidea}[Arc Length]\label{idea:arclength}%
Let $f$ be differentiable on an open interval containing $[a,b]$, where $\fp$ is also continuous on $[a,b]$. Then the arc length of $f$ from $x=a$ to $x=b$ is
\index{integration!arc length}\index{arc length}
\[L = \int_a^b \sqrt{1+[\fp(x)]^2}\dd x.\]
\end{keyidea}

\youtubeVideo{PwmCZAWeRNE}{Arc Length}

As the integrand contains a square root, it is often difficult to use the formula in \autoref{idea:arclength} to find the length exactly. When exact answers are difficult to come by, we resort to using numerical methods of approximating definite integrals. The following examples will demonstrate this.

\begin{example}[Finding arc length]\label{ex_arc1}%
Find the arc length of $f(x) = x^{3/2}$ from $x=0$ to $x=4$.
%
\mtable{A graph of $f(x) = x^{3/2}$ from \autoref{ex_arc1}.}{fig:arc1}{\begin{tikzpicture}[alt={Blue y = x^(3/2) curve on 0 < x < 4 rises steeply, reaching 8 at x = 4.}]
\begin{axis}[width=1.16\marginparwidth,tick label style={font=\scriptsize},
axis y line=middle,axis x line=middle,name=myplot,axis on top,
ymin=-.2,ymax=8.5,xmin=-.1,xmax=4.5,]
\addplot [draw={\colorone},smooth,thick,domain=0:4,samples=50] {x^(3/2)};
%\addplot [draw={\colorone},smooth,thick,domain=0:4,samples=50] coordinates {(0,0)(0.1,0.03162)(0.2,0.08944)(0.3,0.1643)(0.4,0.253)(0.5,0.3536)(0.6,0.4648)(0.7,0.5857)(0.8,0.7155)(0.9,0.8538)(1.,1.)(1.3,1.482)(1.6,2.024)(1.9,2.619)(2.2,3.263)(2.5,3.953)(2.8,4.685)(3.1,5.458)(3.4,6.269)(3.7,7.117)(4.,8.)};
\end{axis}
\node [right] at (myplot.right of origin) {\scriptsize $x$};
\node [above] at (myplot.above origin) {\scriptsize $y$};
\end{tikzpicture}}
%
\solution
A graph of $f$ is given in \autoref{fig:arc1}.
We begin by finding $\fp(x)= \frac32x^{1/2}$. Using the formula, we find the arc length $L$ as
\begin{align*}
	L
	&=	\int_0^4 \sqrt{1+\left(\frac32x^{1/2}\right)^2}\dd x \\
	&=	\int_0^4 \sqrt{1+\frac94x}\ \dd x \\
	&= 	\int_0^4 \left(1+\frac94x\right)^{1/2}\dd x \\
	&=  \frac23\cdot\frac49\left(1+\frac94x\right)^{3/2}\Big|_0^4 \\
	&=	\frac{8}{27}\left(10^{3/2}-1\right)\text{units}. % \approx 9.07 
\end{align*}
\end{example}

\begin{example}[Finding arc length]\label{ex_arc2}%
Find the arc length of $\ds f(x) =\frac18x^2-\ln x$ from $x=1$ to $x=2$.
\solution
A graph of $f$ is given in \autoref{fig:arc2}; the portion of the curve measured in this problem is in bold.
\mtable{A graph of $f(x) =\frac18x^2-\ln x$ from \autoref{ex_arc2}.}{fig:arc2}{\begin{tikzpicture}[alt={Blue y = x²/8 − ln x on 0 < x < 3 dives from 1, bottoms near (2, -0.2), then turns upward.}]
\begin{axis}[width=\marginparwidth,tick label style={font=\scriptsize},
axis y line=middle,axis x line=middle,name=myplot,axis on top,
ymin=-.26,ymax=1.2,xmin=-.1,xmax=3.1]
\addplot [draw={\colorone},smooth,thin,domain=.2:3] {x*x/8-ln(x)};
\addplot [draw={\colorone},smooth,very thick,domain=1:2] {x*x/8-ln(x)};
\end{axis}
\node [right] at (myplot.right of origin) {\scriptsize $x$};
\node [above] at (myplot.above origin) {\scriptsize $y$};
\end{tikzpicture}}
This function was chosen specifically because the resulting integral can be evaluated exactly. We begin by finding $\fp(x) = x/4-1/x$. The arc length is 
\begin{align*}
	L
	&=  \int_1^2 \sqrt{1+ \left(\frac x4-\frac1x\right)^2}\dd x \\
	&= 	\int_1^2 \sqrt{ 1 + \frac{x^2}{16} -\frac12 + \frac1{x^2} } \dd x \\
	&=	\int_1^2 \sqrt{ \frac{x^2}{16} +\frac12 + \frac1{x^2} } \dd x \\
	&=	\int_1^2	\sqrt{ \left(\frac x4 + \frac1x\right)^2 }\dd x \\
	&= \int_1^2 \left(\frac x4 + \frac1x\right) \dd x \\
	&=  \left(\frac{x^2}8 + \ln x\right)\Bigg|_1^2\\
	&=	\frac38+\ln 2\ \text{units}. % \approx 1.07
\end{align*}
\end{example}

The previous examples found the arc length exactly through careful choice of the functions. In general, exact answers are much more difficult to come by and numerical approximations are necessary.

\begin{example}[Approximating arc length numerically]\label{ex_arc3}%
Find the length of the sine curve from $x=0$ to $x=\pi$.
\solution
This is somewhat of a mathematical curiosity; \autoref{ex_ftc4} showed that the area under one ``hump'' of the sine curve is 2 square units; now we are measuring its arc length.

\mtable{A table of values of $y=\sqrt{1+\cos^2x}$ to evaluate a definite integral in \autoref{ex_arc3}.}{fig:arc3}{%
\[\begin{array}{cc}
x & \sqrt{1+\cos^2x} \\ \midrule
 0 & \sqrt{2} \\
 \pi/4 & \sqrt{3/2} \\
 \pi/2 & 1 \\
 3 \pi/4 & \sqrt{3/2} \\
 \pi  & \sqrt{2} \\
\end{array}\]}

The setup is straightforward: $f(x) = \sin x$ and $\fp(x) = \cos x$. Thus 
\[L = \int_0^\pi \sqrt{1+\cos^2x}\dd x.\]
This integral \emph{cannot} be evaluated in terms of elementary functions so we will approximate it with Simpson's Method with $n=4$. \autoref{fig:arc3} gives the integrand evaluated at 5 evenly spaced points in $[0,\pi]$. Simpson's Rule then states that 
\begin{align*}
	\int_0^\pi \sqrt{1+\cos^2x}\dd x
	&\approx	\frac{\pi-0}{4\cdot 3}
	\left(\sqrt{2}+4\sqrt{3/2}+2(1)+4\sqrt{3/2}+\sqrt{2}\right) \\
	&\approx3.82918.
\end{align*}
Using a computer with $n=100$ the approximation is $L\approx 3.8202$; our approximation with $n=4$ is quite good.  Our approximation of $3.79$ from the beginning of this section isn't as close.
\end{example}

\subsection{Surface Area of Solids of Revolution}

\mtable{Establishing the formula for surface area.}{fig:surface_intro}{%
\myincludeasythree{
3Droll=126.0060482236849,
3Dortho=0.004875946324318647,
3Dc2c=0.44462037086486816 0.39410829544067383 0.8043577671051025,
3Dcoo=67.21983337402344 3.5258262157440186 -29.566415786743164,
3Droo=150.00000160608386}{\marginparwidth}{Blue arc y = f(x) between x = a and x = b; short red chord shows one sample slice. Revolving that slice about the x-axis gives a purple band on a blue tubular surface.}{figures/figarc4_b_3D}
% todo Tim should this be 3D at all?
\\(a)\\
\myincludeasythree{
3Droll=126.0060482236849,
3Dortho=0.004875946324318647,
3Dc2c=0.44462037086486816 0.39410829544067383 0.8043577671051025,
3Dcoo=67.21983337402344 3.5258262157440186 -29.566415786743164,
3Droo=150.00000160608386}{\marginparwidth}{Blue band on blue tube: frustum formed when the red chord of the slice is revolved about the x-axis}{figures/figarc4_3D}
\\(b)}

We have already seen how a curve $y=f(x)$ on $[a,b]$ can be revolved around an axis to form a solid. Instead of computing its volume, we now consider its surface area.

We begin as we have in the previous sections: we partition the interval $[a,b]$ with $n$ subintervals, where the $i\,^{\text{th}}$ subinterval is $[x_i,x_{i+1}]$. On each subinterval, we can approximate the curve $y=f(x)$ with a straight line that connects $f(x_i)$ and $f(x_{i+1})$ as shown in \autoref{fig:surface_intro}(a). Revolving this line segment about the $x$-axis creates part of a cone (called a \emph{frustum} of a cone) as shown in \autoref{fig:surface_intro}(b). The surface area of a frustum of a cone is
\[A=2\pi r_{\text{avg}} L,\]
where $r_{\text{avg}}$ is the average of $R_1$ and $R_2$.  The length is given by $L$; we use the material just covered by arc length to state that
\[L\approx \sqrt{1+[\fp(c_i)]^2}\Delta x_i\]
for some $c_i$ in the $i^{\text{th}}$ subinterval. The radii are just the function evaluated at the endpoints of the interval: $f(x_{i-1})$ and $f(x_i)$. Thus the surface area of this sample frustum of the cone is approximately 
\[2\pi\frac{f(x_{i-1})+f(x_i)}2\sqrt{1+[\fp(c_i)]^2}\Delta x_i.\]

Since $f$ is a continuous function, the Intermediate Value Theorem states there is some $d_i$ in $[x_{i-1},x_i]$ such that $f(d_i)=\dfrac{f(x_{i-1})+f(x_i)}2$; we can use this to rewrite the above equation as
\[2\pi f(d_i)\sqrt{1+[\fp(c_i)]^2}\Delta x_i.\]
Summing over all the subintervals we get the total surface area to be approximately 
\[\text{Surface Area}\approx \sum_{i=1}^n 2\pi f(d_i)\sqrt{1+[\fp(c_i)]^2}\Delta x_i,\]
which is almost a Riemann Sum (we would need $d_i=c_i$ to remove the ``almost''). Taking the limit as the subinterval lengths go to zero gives us the exact surface area, given in the upcoming Key Idea.

If instead we revolve $y=f(x)$ about the $y$-axis, the radii of the resulting frustum are $x_{i-1}$ and $x_i$; their average value is simply the midpoint of the interval. In the limit, this midpoint is just $x$. This gives the second part of \autoref{idea:surface_area}.

% todo Tim does this (and prev KI) only need f' continuous on (a,b)?
\begin{keyidea}[Surface Area of a Solid of Revolution]\label{idea:surface_area}%
Let $f$ be differentiable on an open interval containing $[a,b]$ where $\fp$ is also continuous on $[a,b]$. \index{integration!surface area}\index{surface area!solid of revolution}
\begin{enumerate}
	\item	The surface area of the solid formed by revolving the graph of $y=f(x)$, where $f(x)\geq0$, about the $x$-axis is
	\[\text{Surface Area} = 2\pi\int_a^b f(x)\sqrt{1+[\fp(x)]^2}\dd x.\]
	\item	The surface area of the solid formed by revolving the graph of $y=f(x)$ about the $y$-axis, where $a,b\geq0$, is
	\[\text{Surface Area} = 2\pi\int_a^b x\sqrt{1+[\fp(x)]^2}\dd x.\]
\end{enumerate}
\end{keyidea}

\begin{example}[Finding surface area of a solid of revolution]\label{ex_sa1}%
Find the surface area of the solid formed by revolving $y=\sin x$ on $[0,\pi]$ around the $x$-axis, as shown in \autoref{fig:sa1}.
\solution
The setup turns out to be easier than the resulting integral. Using \autoref{idea:surface_area}, we have the surface area $SA$ is:
\mtable{Revolving $y=\sin x$ on $[0,\pi]$ about the $x$-axis.}{fig:sa1}{%
\myincludeasythree{
3Droll=121.00092658103458,
3Dortho=0.0048407199792563915,
3Dc2c=0.2899017035961151 0.17925968766212463 0.9401186108589172,
3Dcoo=72.58285522460938 2.891160011291504 -26.545438766479492,
3Droo=149.99998703790175}{\marginparwidth}{Blue sinus-shaped surface: y = sin x (0 < x < pi) revolved about the x-axis, axes shown.}{figures/figsa1_3D}}
\begin{align*}
	\text{SA}
	&=	2\pi\int_0^\pi \sin x\sqrt{1+\cos^2x}\dd x \\
	&=	-2\pi\int_1^{-1}\sqrt{1+u^2}\dd u \qquad\qquad \text{substitute $u=\cos x$} \\
	&=	2\pi\int_{-\pi/4}^{\pi/4}\sec^3\theta\dd\theta \qquad\qquad \text{substitute $u=\tan\theta$} \\
	&=  \left.\pi\left(\sec\theta\tan\theta+\ln\abs{\sec\theta+\tan\theta}\right)\right|_{-\pi/4}^{\pi/4} \qquad \text{by \autoref{ex_trigint6}} \\
	&=  \pi\left(\sqrt2+\ln\left(\sqrt2+1\right)
	-\left(-\sqrt2+\ln\left(\sqrt2-1\right)\right)\right) \\
	&=	\pi\left(2\sqrt2+\ln\left(\frac{\sqrt2+1}{\sqrt2-1}\right)\right) \\
	&=	2\pi\left(\sqrt2+\ln\left(\sqrt2+1\right)\right)\ \text{units}^2 \qquad\qquad\text{rationalize the denominator.}
\end{align*}

It is interesting to see that the surface area of a solid, whose shape is defined by a trigonometric function, involves both a square root and a natural logarithm.
\end{example}

\mtable{The solids used in \autoref{ex_sa2}.}{fig:sa2}{%
\myincludeasythree{
3Droll=138.26595091591264,
3Dortho=0.0048407199792563915,
3Dc2c=0.16831636428833008 0.19470007717609406 0.966313362121582,
3Dcoo=81.32470703125 3.403531074523926 -28.20997428894043,
3Droo=149.9999886313645}{\marginparwidth}{Blue horn-like surface: y = x^2 (0 < x < 1) rotated about x-axis, tip at (0,0,0) flaring to x = 1.}{figures/figsa2a_3D}
\\(a)\\[15pt]
\myincludeasythree{
3Droll=138.26595091591264,
3Dortho=0.0048407199792563915,
3Dc2c=0.16831637918949127 0.19470007717609406 0.966313362121582,
3Dcoo=0.013853734359145164 79.42704010009766 -29.364774703979492,
3Droo=149.99998895240248}{\marginparwidth}{Blue parabolic bowl: y = x^2 (0 < x < 1) revolved about the y-axis; vertex at origin, rim near y = 1.}{figures/figsa2b_3D}\\
(b)}

\begin{example}[Finding surface area of a solid of revolution]\label{ex_sa2}%
Find the surface area of the solid formed by revolving the curve $y=x^2$ on $[0,1]$ about:\\
\null\hfill 1.\quad the $x$-axis\hfill2.\quad the $y$-axis.\hfill\null
\solution
\begin{enumerate}
	\item		The solid formed by revolving $y=x^2$ around the $x$-axis is graphed in \autoref{fig:sa2}(a). Like the integral in \autoref{ex_sa1}, this integral is easier to setup than to actually integrate. While it is possible to use a trigonometric substitution to evaluate this integral, it is significantly more difficult than a solution employing the hyperbolic sine:
\begin{align*}
	SA &= 2\pi\int_0^1 x^2\sqrt{1+(2x)^2}\dd x. \\
%	&= 2\pi\int_0^{\tan^{-1}2} \frac18\tan^2\theta\sec^3\theta\ d\theta\qquad\text{substitute $x=\frac12\tan\theta$} \\
%	&= \frac\pi4\int_0^{\tan^{-1}2}\sec^5\theta-\sec^3\theta\ d\theta \\
	&= \left.\frac{\pi}{32}\left(2(8x^3+x)\sqrt{1+4x^2}-\sinh^{-1}(2x)\right)\right|_0^1\\
	&=\frac{\pi}{32}\left(18\sqrt{5}-\sinh^{-1}2\right)\ \text{units}^2.
\end{align*}

	\item	 Since we are revolving around the $y$-axis, the ``radius'' of the solid is not $f(x)$ but rather $x$. Thus the integral to compute the surface area is:
\begin{align*}
	SA &= 2\pi\int_0^1x\sqrt{1+(2x)^2}\dd x \\
	&=	\frac{\pi}4\int_1^5 \sqrt{u}\dd u \qquad\text{substitute $u=1+4x^2$} \\
	&= \left.\frac{\pi}{4}\frac23 u^{3/2}\right|_1^5\\
	&= \frac{\pi}6\left(5\sqrt{5}-1\right)\ \text{units}^2.
\end{align*}
 The solid formed by revolving $y=x^2$ about the $y$-axis is graphed in \autoref{fig:sa2} (b).
\end{enumerate}
\end{example}

Our final example is a famous mathematical ``paradox.''

\mtable{A graph of Gabriel's Horn.}{fig:gabriel}{%
\myincludeasythree{
3Droll=107.09033109095482,
3Dortho=0.0048407199792563915,
3Dc2c=0.4647676646709442 0.1565970778465271 0.8714748024940491,
3Dcoo=72.58285522460938 2.8911631107330322 -26.545440673828125,
3Droo=149.99999776925813}{\marginparwidth}{Infinite blue trumpet (Gabriel’s Horn): surface from rotating y = 1 / x, 1 < x < infty, about the x-axis.}{figures/figgabriel_3D}}
\begin{example}[The surface area and volume of Gabriel's Horn]\label{ex_gabriel}%
Consider the solid formed by revolving $y=1/x$ about the $x$-axis on $[1,\infty)$. Find the volume and surface area of this solid. (This shape, as graphed in \autoref{fig:gabriel}, is known as ``Gabriel's Horn'' since it looks like a very long horn that only a supernatural person, such as an angel, could play.)\index{Gabriel's Horn}
\solution
To compute the volume it is natural to use the Disk Method. We have:
\begin{align*}
V &= \pi\int_1^\infty \frac{1}{x^2}\dd x \\
	&= \lim_{b\to\infty}\pi\int_1^b\frac{1}{x^2}\dd x \\
	&=	\lim_{b\to\infty} \left.\pi\left(\frac{-1}{x}\right)\right|_1^b \\
	&= \lim_{b\to\infty} \pi\left(1-\frac1b\right) \\
	&= \pi \ \text{units}^3.
\end{align*}
Gabriel's Horn has a finite volume of $\pi$ cubic units. Since we have already seen that regions with infinite length can have a finite area, this is not too difficult to accept.

We now consider its surface area. The integral is straightforward to setup:
\begin{align*}
SA &= 2\pi\int_1^\infty \frac{1}{x}\sqrt{1+1/x^4}\dd x.
\intertext{Integrating this expression is not trivial. We can, however, compare it to other improper integrals. Since \intertextmath{1< \sqrt{1+1/x^4}} on \intertextmath{[1,\infty)}, we can state that}
2\pi\int_1^\infty \frac{1}{x}\dd x &<2\pi\int_1^\infty \frac{1}{x}\sqrt{1+1/x^4}\dd x .
\end{align*}
By \autoref{idea:impint1}, the improper integral on the left diverges. Since the integral on the right is larger, we conclude it also diverges, meaning Gabriel's Horn has infinite surface area.

Hence the ``paradox'': we can fill Gabriel's Horn with a finite amount of paint, but since it has infinite surface area, we can never paint it.

Somehow this paradox is striking when we think about it in terms of volume and area. However, we have seen a similar paradox before, as referenced above. We know that the area under the curve $y=1/x^2$ on $[1,\infty)$ is finite, yet the shape has an infinite perimeter. Strange things can occur when we deal with the infinite.
\end{example}

%A standard equation from physics is ``Work = force $\times$ distance'', when the force applied is constant. In the next section we learn how to compute work when the force applied is variable.

\printexercises{exercises/07-04-exercises}
